% !TeX program = pdflatex
% Reasoning/memory separation — a short theory note.
% 分類紀律（Codex [197]/[198]）：每個中心句必須是
%   Definition | Proposition/Theorem + proof | Imported result + exact assumptions | Conjecture
\documentclass[11pt]{article}
\usepackage[margin=1in]{geometry}
\usepackage{amsmath,amssymb,amsthm}
\usepackage{hyperref}

\theoremstyle{definition}
\newtheorem{definition}{Definition}
\newtheorem{assumption}{Assumption}
\newtheorem{conjecture}{Conjecture}
\theoremstyle{remark}
\newtheorem{remark}{Remark}
\theoremstyle{plain}
\newtheorem{lemma}{Lemma}
\newtheorem{theorem}{Theorem}

\newcommand{\Dom}{\mathcal{D}}
\newcommand{\Wr}{\mathcal{W}}
\newcommand{\Qu}{\mathcal{Q}}
\newcommand{\Ky}{\mathcal{K}}
\newcommand{\Va}{\mathcal{V}}
\newcommand{\bt}{\bot}
\newcommand{\cw}{C_w}
\newcommand{\cr}{C_r}

\title{Reasoning/Memory Separation in Language Models:\\
an address-factorization theorem for a restricted interface}
\author{(draft)}
\date{}

\begin{document}
\maketitle

\begin{abstract}
We give a minimal formal setting in which ``can knowledge be separated from a
language model into an external, addressable memory?'' becomes a precise
question about the \emph{addressing interface} rather than about the model.
We prove (i) an elementary state-independence lemma, and (ii) an
address-factorization theorem: for one-shot key-addressed systems over a
\emph{fixed public schema}, with no store-bypassing channel and zero error
required over all finite semantic states, universal completeness and
non-mis-delivery hold \emph{iff} the interface realises the schema's naming
map on certified inputs. The theorem is a representation result for that
restricted class. It does \emph{not} claim that the schema, the value codec,
or natural-language addressing can be learned, and it makes no claim about
approximate correctness.
\end{abstract}

\section{Motivation (prose)}

A recurring engineering picture is: a small model that reasons, plus an
external store that remembers. The picture is normally defended empirically.
Our question is prior: \emph{under what conditions can such a system be
correct at all?} We fix a vocabulary, prove the one step that is genuinely
forced, and reduce the remainder to a single named property of the interface.

\section{Task specification and semantics}

\begin{definition}[Task specification]\label{def:task}
A \emph{task} fixes, publicly and in advance: a referent domain $\Dom$, write
utterances $\Wr$, query utterances $\Qu$, a key space $\Ky$, a value space
$\Va$ with $|\Va|\ge 2$, partial semantics
$s_w:\Wr\rightharpoonup\Dom$, $s_q:\Qu\rightharpoonup\Dom$, and an
\emph{injective public naming map} $\iota:\Dom\hookrightarrow\Ky$.
\end{definition}

\noindent$\iota$ is a schema (an oracle specification), \emph{not} a learned
canonicaliser. What is at issue below is whether an interface \emph{realises}
$\iota$.

\begin{definition}[Semantic state and legal store]\label{def:state}
A \emph{semantic state} is a finite partial map $S:\Dom\rightharpoonup\Va$.
Its \emph{legal store} is the partial map $M_S:\Ky\rightharpoonup\Va$ with
$M_S(\iota(d))=S(d)$ for $d\in\mathrm{dom}(S)$ and undefined elsewhere.
Memory is manipulated only by the typed operations $\mathrm{put}(M,k,v)$ and
$\mathrm{get}(M,k)\in\Va\cup\{\bt\}$, where $\mathrm{get}(M,k)=\bt$ exactly
when $k\notin\mathrm{dom}(M)$.
\end{definition}

\begin{definition}[Task answer]\label{def:answer}
For a semantic state $S$ and query $q$,
\[
A(S,q)=\begin{cases} S(s_q(q)) & s_q(q)\in\mathrm{dom}(S),\\
\bt & \text{otherwise.}\end{cases}
\]
\end{definition}

\noindent $A$ is given by the task semantics alone; it does not mention any
interface. This is what keeps Lemma~\ref{lem:state} non-circular.

\begin{definition}[Addressing interface]\label{def:interface}
An \emph{addressing interface} is a pair of total maps
$\cw:\Wr\to\Ky\cup\{\bt\}$ and $\cr:\Qu\to\Ky\cup\{\bt\}$.
\end{definition}

\begin{definition}[Certificate relation]\label{def:cert}
The task additionally fixes externally given, decidable relations
$\mathsf{cert}_w\subseteq\Wr\times\Ky$ and
$\mathsf{cert}_r\subseteq\Qu\times\Ky$. An interface is \emph{certified} if
$\cw(w)=k\neq\bt$ only when $\mathsf{cert}_w(w,k)$ holds, and likewise for
$\cr$; when no certified key exists the interface outputs $\bt$.
\end{definition}

\begin{assumption}[Certificate soundness]\label{as:sound}
$\mathsf{cert}_w(w,k)\Rightarrow k=\iota(s_w(w))$, and
$\mathsf{cert}_r(q,k)\Rightarrow k=\iota(s_q(q))$, whenever the semantics are
defined.
\end{assumption}

\begin{assumption}[In-scope coverage]\label{as:cov}
For every in-scope $w$ (resp.\ $q$) with $s_w(w)$ (resp.\ $s_q(q)$) defined,
some $k$ with $\mathsf{cert}_w(w,k)$ (resp.\ $\mathsf{cert}_r(q,k)$) exists.
\end{assumption}

\noindent Assumptions~\ref{as:sound}--\ref{as:cov} are \emph{premises}, not
results. In an open world nothing above supplies them; ``the referent is
genuinely unknown'' is not established by this note.

\begin{definition}[One-shot key-addressed system]\label{def:system}
A system is \emph{one-shot key-addressed} if on write $w$ with value $v$ it
performs $\mathrm{put}(M,\cw(w),v)$ (a no-op if $\cw(w)=\bt$), on query $q$ it
outputs $\mathrm{get}(M,\cr(q))$ if $\cr(q)\neq\bt$ and $\bt$ otherwise, and
it has \emph{no channel to $M$} other than that single $\mathrm{get}$
(\emph{no store-bypassing side channel}).
\end{definition}

\section{What is forced}

\begin{lemma}[State-independence]\label{lem:state}
Let $F:\Qu\to\Va\cup\{\bt\}$ be any function that does not depend on the
runtime memory. If there are semantic states $S_0,S_1$ and a query $q$ with
$A(S_0,q)\neq A(S_1,q)$, then $F$ is not exact-correct on both.
\end{lemma}

\begin{proof}
$F(q)$ is a single element of $\Va\cup\{\bt\}$ and cannot equal two distinct
values.
\end{proof}

\begin{remark}[Reading, deliberately narrow]
Lemma~\ref{lem:state} shows only that exact correctness across semantic states
requires \emph{some} channel through which the runtime state influences the
output. It does \emph{not} entail that the channel be a key--value store, a
vector database, or a non-neural guard, and it says nothing about approximate
correctness. It is elementary; we state it because it is the only step of the
usual argument for externalised memory that is actually forced.
\end{remark}

\section{Certified canonical addressability}

\begin{definition}[CCA]\label{def:cca}
A certified interface $(\cw,\cr)$ has \emph{Certified Canonical
Addressability} with respect to $\iota$ if, on every in-scope certified input,
\[
\cw(w)=\iota(s_w(w)),\qquad \cr(q)=\iota(s_q(q)),
\]
and $\bt$ is output exactly when no certified key exists.
\end{definition}

\begin{remark}[Two readable consequences]\label{rem:readable}
CCA implies (i) \emph{same-referent agreement}: if $s_w(w)=s_q(q)$ then
$\cw(w)=\cr(q)\neq\bt$; and (ii) \emph{live-referent separation}: distinct
referents receive distinct keys. Clause (ii) is not redundant --- agreement
alone is satisfied by the degenerate interface mapping every utterance to one
fixed key, which is perfectly consistent and useless.
\end{remark}

\noindent We do \emph{not} let the interface decide which key is canonical;
otherwise the theorem below degenerates into renaming.

\section{Address-factorization theorem}

\begin{definition}[Universal soundness and completeness]\label{def:usc}
Fix a task specification. A one-shot key-addressed system is \emph{universally
sound and complete} if for every finite semantic state $S$, every admissible
history of atomic $\mathrm{put}/\mathrm{get}$ consistent with $M_S$, and every
in-scope $q$ with $s_q(q)$ defined, the system outputs exactly $A(S,q)$.
\end{definition}

\begin{theorem}[Universal-store address factorization]\label{thm:main}
Under Assumptions~\ref{as:sound}--\ref{as:cov}, a one-shot key-addressed
system over a fixed public schema, with no store-bypassing channel, is
universally sound and complete \emph{iff} its interface has CCA
(Definition~\ref{def:cca}).
\end{theorem}

\begin{proof}
($\Leftarrow$) Assume CCA. Let $S$ be any semantic state and $q$ in scope.
If $s_q(q)\in\mathrm{dom}(S)$ then by Assumption~\ref{as:cov} a certified key
exists, and by CCA $\cr(q)=\iota(s_q(q))$, so the system outputs
$M_S(\iota(s_q(q)))=S(s_q(q))=A(S,q)$. If $s_q(q)\notin\mathrm{dom}(S)$ then
$\iota(s_q(q))\notin\mathrm{dom}(M_S)$, so $\mathrm{get}$ returns $\bt=A(S,q)$.

($\Rightarrow$) Suppose CCA fails, i.e.\ there is an in-scope input on which
the interface disagrees with $\iota$. Write $d$ for its referent.

\emph{Case 1: $\cr(q)=\iota(d')$ with $d'\neq d$.} Since $|\Va|\ge2$ choose
$v\neq v'$ and take $S$ with $\mathrm{dom}(S)=\{d,d'\}$, $S(d)=v$,
$S(d')=v'$. The system outputs $M_S(\iota(d'))=v'\neq v=A(S,q)$: a
mis-delivery of a value that is genuinely present.

\emph{Case 2: $\cr(q)\notin\iota(\Dom)$, or $\cr(q)=\bt$ for a covered $q$.}
Take any $S$ with $d\in\mathrm{dom}(S)$. Then the system outputs $\bt$ while
$A(S,q)=S(d)\neq\bt$: incompleteness.

\emph{Case 3: $\cw(w)\neq\iota(s_w(w))$ for a covered $w$.} The written value
is absent from $\iota(s_w(w))$, so a later query for $s_w(w)$ reduces to
Case~2; and if $\cw(w)=\iota(d'')$ for some $d''\neq s_w(w)$, a query for
$d''$ reduces to Case~1.

In each case universal soundness and completeness fails.
\end{proof}

\begin{remark}[Responsibility boundary]\label{rem:boundary}
Case~1 of the proof is the operative one for safety. If a query for $d$ is
routed to a \emph{live} key of $d'\neq d$, exact membership testing and
$\mathrm{get}$ both behave correctly and still return the wrong value. Hence
an exact-membership guard can certify that a \emph{key} is absent; it cannot
certify that a present key is \emph{the intended referent}, and it cannot
repair an upstream identity error.
\end{remark}

\begin{remark}[Private stores: agreement up to relabeling]\label{rem:relabel}
Theorem~\ref{thm:main} is stated for a \emph{fixed public schema}: legal
stores are defined through $\iota$, which is what makes memory shareable and
auditable across systems. If instead the store is private and populated only
by this system's own writes, then any fixed injection $\pi$ with
$\cw=\cr=\pi\circ\iota$ also yields universal correctness, and CCA can be
weakened to \emph{agreement up to a fixed injection}. The public-schema
reading is the one that supports interoperability; we record the weaker
private-store variant explicitly rather than leave the distinction implicit.
\end{remark}

\paragraph{Scope.} Theorem~\ref{thm:main} is a representation/refinement
result for a restricted class: fixed public schema, one-shot key-addressed,
no bypass, zero error. It proves interface correctness \emph{relative to a
schema}. It does \emph{not} prove that the schema or natural-language
semantics can be learned, that a value codec can be learned, or that a
reasoning executor can consume latent values.

\section{Imported results (to be filled with exact assumptions)}

\emph{Placeholder.} Each item will be stated with the assumptions under which
it holds: precision and attention assumptions for transformer circuit upper
bounds; arbitrary-precision activations for Turing-completeness
constructions; quantisation premises for parameter-capacity counting;
conformal-risk (not semantic-identity) guarantees for retrieval systems.
Architectures such as differentiable read/write memories, and empirical
external-memory results, are adjacent but do \emph{not} establish CCA or
universal non-mis-delivery.

\paragraph{Prohibited inference (recorded to prevent it).} A
$\mathsf{TC}^0$ \emph{upper} bound for fixed-depth transformers combined with
$\mathsf{NC}^1$-completeness of a word problem does \emph{not} give an
unconditional lower bound: the required class separation is not known.

\section{Conjectures}

\begin{conjecture}\label{conj:scale}
Whether the same-referent agreement consequence of CCA
(Remark~\ref{rem:readable}(i)) arises at scale from unsupervised training is
open. We claim nothing in either direction.
\end{conjecture}

\section{What this note does not claim}

That a separable system is achievable for natural language; that it is
impossible; that latent delivery is superior to placing the value in context;
or that any measured failure of a particular mechanism bears on
Definition~\ref{def:cca} beyond instantiating it.

\end{document}
